\chapter{Inlining and Specialization}\label{ch:inlining-specialization}

A procedure call looks cheap in source code but costs real work at run time.
The caller saves registers, passes arguments, jumps to a new address, builds a frame, and later tears it down.
Inlining removes that boundary by splicing the body of a callee into its caller.
Specialization goes one step further and builds a dedicated copy of a procedure for a particular constant, type, or shape.
Both transformations trade code size for speed, and both unlock other optimizations that cannot see through a call.

\section{Why Calls Block Optimization}\label{sec:ch23-why-calls-block}

Consider a tiny helper that clamps a value into a range.
The caller passes two bounds and a candidate, and the callee does two comparisons.
Written with a call, constant propagation stops at the call boundary because the optimizer has no facts about the callee body at the call site.
Written inline, the comparisons sit next to the surrounding constants and fold away.
That pattern repeats for loop bounds, null checks, type tags, and array lengths.
The optimization that matters is often not the eliminated jump but the cleanup that becomes visible afterward.
\term{Inlining} replaces a call site with a renamed copy of the callee body.
\term{Specialization} creates a new procedure variant tuned to a particular argument property and redirects selected call sites to it.
\term{Direct call} names its target statically, while \term{indirect call} selects the target at run time through a pointer or method table.
Direct calls are easy to inline.
Indirect calls need a guard, a profile, or a devirtualization proof first.

The following three-address fragment shows the lost opportunity.
\begin{irlisting}
; caller before inlining
t1 = call clamp(x, 0, 100)
t2 = add t1, 1
; callee, defined elsewhere
clamp(a, lo, hi):
  t3 = cmp_lt a, lo
  t4 = select t3, lo, a
  t5 = cmp_gt t4, hi
  t6 = select t5, hi, t4
  ret t6
\end{irlisting}
After inlining, the two bounds are the constants \code{0} and \code{100} sitting next to the comparisons.
A later pass folds, simplifies selects, and may prove that one branch is dead for this call site.
Without inlining, none of that happens.
The same effect appears with wrapper functions, accessor methods, and small numeric kernels.

Inlining also changes the shape of the control-flow graph.
A call node becomes straight-line code, a fork, or a small diamond, depending on the callee.
That new shape can merge with the caller loops and branches.
Loop-invariant code motion, global value numbering, and dead-code elimination all profit from the larger scope.
The price is code growth and possibly worse instruction-cache behavior.
A good inliner therefore weighs the expected cleanup against the size cost of each call site.

\section{The Inlining Decision}\label{sec:ch23-decision}

An inliner needs three pieces: legality, profitability, and mechanics.
Legality asks whether the transformation preserves meaning.
Profitability asks whether it is worth doing.
Mechanics performs the splice without breaking names, control flow, or exception behavior.
The three pieces interact because a profitable-looking site may be illegal without extra guards.
Figure~\ref{fig:ch23-inline-scope} shows how inlining widens the optimizer view from two isolated regions into one joint region.
Table~\ref{tab:ch23-inline-tradeoffs} summarizes the main benefits and the corresponding risks.

\begin{figure}[htbp]
\centering
\begin{tikzpicture}[node distance=1.1cm and 1.6cm]
\node[basic block] (caller) {caller: t = call f(a)};
\node[basic block, right=of caller] (callee) {f: body; ret r};
\node[basic block, below=of caller] (merged) {caller: body[a/x]; use r};
\draw[cfg edge] (caller) -- (callee);
\draw[cfg edge] (caller) -- (merged);
\end{tikzpicture}
\caption{Inlining merges the callee body into the caller, so later passes see one region instead of two isolated ones.}
\label{fig:ch23-inline-scope}
\end{figure}

\begin{table}[htbp]
\caption{Benefits of inlining and specialization with their main risks.}
\label{tab:ch23-inline-tradeoffs}
\begin{tabular}{lll}
\toprule
Benefit & Mechanism & Risk \\
\midrule
Eliminated call overhead & No jump or frame setup & Larger code size \\
Exposed constants & Arguments become local facts & Extra register pressure \\
Better branch folding & Callee tests meet caller values & More hot-path branches \\
Type refinement & Clone assumes one case & Guard cost on mismatch \\
Cross-call code motion & One region for hoisting & Slower compile time \\
\bottomrule
\end{tabular}
\end{table}

Legality is straightforward for direct calls to nonrecursive functions without tricky control flow.
The compiler renames the callee locals so they do not collide with caller names.
It maps actual arguments to fresh temporaries, splices the body in place of the call, and wires each \code{ret} to the continuation.
Recursion needs a depth bound or a partial strategy that inlines a few levels and leaves a residual call.
Functions with multiple returns need one join point per inlined site or a duplicated tail.
Language features such as exceptions, coroutines, and dynamic scoping add edges that the splice must preserve.
A conservative inliner refuses sites it cannot model exactly.
\begin{ccwarning}[Do not inline what you cannot unwind]
Inlined exception edges, cleanup actions, and early returns must land in the right handler.
If the intermediate representation cannot represent the combined control flow, keep the call.
A wrong landing pad corrupts unwinding tables and exception semantics.
\end{ccwarning}

Profitability is a heuristic because the true payoff depends on passes that have not run yet.
Most inliners combine a size budget with a benefit score.
Small callees get inlined almost always.
Large callees need strong evidence such as constant arguments, a single call site, or a hot profile count.
Call-site context matters: the same callee may be worth inlining in a tight loop and not worth inlining on a cold error path.
Specialization uses the same logic but pays the code-size cost once per variant and shares it across sites.
A \term{call-site benefit} estimates the constants, branches, and loads that become simplifiable after the splice.
A \term{size budget} caps total growth for a function or module.

Specialization handles the case where full inlining is too big but one argument property makes a cheaper version possible.
A generic routine over an element size, a data type tag, or a stride gains a fast path when that property is fixed.
The compiler clones the procedure, replaces the parameter with the known value inside the clone, runs the normal simplifiers on the clone, and redirects matching call sites.
Unmatched sites keep calling the generic version.
Profile data often picks the specialization key, for example the most common receiver type at a virtual call.
The guard that checks the key must be cheaper than the work it saves.

\section{Cost-Guided Inlining Algorithm}\label{sec:ch23-algorithm}

Algorithm~\ref{alg:ch23-inline} states the standard greedy worklist approach.
It scores every eligible direct call site, inlines the best one within budget, and repeats because new sites and new constants appear.
The scoring function is deliberately simple so it runs on every module without profile data.
A profile-guided compiler adds hotness multipliers on top of the same skeleton.
The algorithm assumes three-address or SSA input with explicit calls and a call graph the compiler can update incrementally.
It also assumes the simplifier runs between rounds so that benefit estimates track the current code.

\begin{algorithm}[htbp]
\caption{Greedy cost-guided inlining with a size budget}\label{alg:ch23-inline}
\begin{pseudocode}
\Require call graph $G$, size budget $B$, scoring function $Score$
\Ensure inlined module with updated $G$
\State compute initial size $S$ of each function
\State build worklist $W$ of eligible direct call sites
\While{$W \neq \emptyset$}
\State remove the site $c$ with the highest $Score(c)$ from $W$
\If{$Score(c) \le 0$}
\State \Break
\EndIf
\If{inlining $c$ exceeds budget $B$}
\State \Continue
\EndIf
\State splice a renamed copy of the callee at $c$
\State run local cleanup: fold, propagate, eliminate dead code
\State update sizes $S$, call graph $G$, and affected scores
\State add newly exposed call sites to $W$
\EndWhile
\State emit specialized clones for profitable constant keys
\State redirect matching call sites to the clones
\end{pseudocode}
\end{algorithm}

The score typically subtracts a size penalty from a context bonus.
A common form is a weighted sum of constant arguments, single-call-site status, small callee size, and static estimates of eliminated branches.
Production compilers tune the weights empirically and clamp them so no single factor dominates.
The budget may be global, per caller, or adaptive to optimization level.
A size budget alone is not a termination argument: recursively inlining a wrapper may expose another call without increasing the measured size.
Track recursive expansion depth, deduplicate specialization keys, and impose a finite expansion budget independent of cleanup's size reductions.
Recompute stale scores before using the highest-score stopping rule, and charge specialized clones against the budget as well.
A debug build inlines almost nothing.
A release build with link-time information inlines across modules.
The complexity is $O(C \log C + I \cdot (N + E))$, where $C$ is the number of candidate sites, $I$ is the number of performed inlines, and $N$ and $E$ are the nodes and edges touched by cleanup and graph updates.
The dominant cost is usually the repeated cleanup, not the priority queue.

Specialization reuses the same loop with a different candidate generator.
Each candidate is a pair of a procedure and a specialization key, such as \enquote{second argument equals 4} or \enquote{receiver has a given hidden class}.
The compiler materializes the clone lazily, simplifies it eagerly, and measures whether the clone is smaller or faster than the generic body under that key.
Clones that fail to simplify are discarded before they reach the object file.
This \enquote{simplify then measure} discipline keeps speculative cloning from exploding code size.
% TODO(cite): need a primary reference for profile-guided inlining heuristics.

\section{Worked Example}\label{sec:ch23-worked}

We now inline a small example by hand and track every simplification.
The callee scales and biases its input.
The caller supplies constants at two sites, one of which makes the bias dead.
\begin{irlisting}
; callee
scale_bias(v, s, b):
  t1 = mul v, s
  t2 = add t1, b
  ret t2
; caller before inlining
x1 = call scale_bias(p, 2, 0)
x2 = call scale_bias(q, 4, 7)
\end{irlisting}
Inline the first site by mapping \code{v} to \code{p}, \code{s} to \code{2}, and \code{b} to \code{0} with fresh temporaries.
The spliced code is \code{t1a = mul p, 2} followed by \code{x1 = add t1a, 0}.
Fold the addition of zero to get \code{x1 = mul p, 2}, possibly strength-reduced later to a shift.
No frame, no jump, and no argument shuffling remain.
Inline the second site the same way to get \code{t1b = mul q, 4} and \code{x2 = add t1b, 7}.
Here the addition survives because \code{7} is not an identity, but the multiply by a power of two is still a shift candidate.
The two expansions share no temporaries, so no renaming conflict arises.

Now consider specialization instead of full inlining.
Suppose the real routine is large because it handles strides, checks alignment, and selects among three kernels.
Inlining all of that twice would bloat the caller.
The profile shows that \code{s} is almost always \code{4} at hot sites.
The compiler clones the routine with \code{s} fixed to \code{4} and folds the kernel dispatch to the stride-4 path.
Stride alone does not prove base-pointer alignment; an alignment test remains unless a separate fact establishes the required alignment.
Hot sites call the clone directly.
Cold sites with other strides keep calling the generic version.
Code grows by one compact fast path rather than by two full expansions.
The guard is free here because the call sites pass the constant explicitly; with a dynamic key the compiler would emit a comparison and keep both targets.

We verify the result by counting operations before and after.
Before, each site pays a call plus the two arithmetic operations inside the callee.
After inlining, the first site pays one multiply and the second pays one multiply and one add.
The specialization variant pays one direct call to the lean clone plus the reduced body inside it.
Either way, the arithmetic that depended on the known argument is gone.
That is the signature of a successful inline or specialization: known inputs dissolve neighboring operations.

\section{Implementation Notes}\label{sec:ch23-implementation}

Real inliners live inside a pass pipeline, not in isolation.
They run early to expose constants, again in the middle to catch functions simplified by other passes, and sometimes at link time to cross module boundaries.
LLVM performs inline costing with a modular bonus model and re-runs simplification after each round, while GCC combines early inlining with interprocedural summaries and profile feedback.
Both systems cap growth with budgets that scale with the optimization level.
Hot call sites get larger budgets than cold ones.
Recursive functions get a small unroll-like allowance rather than unbounded expansion.
These details matter because the textbook greedy loop looks optimal but is only a heuristic: inlining under realistic size constraints is NP-hard, so production compilers tune rather than solve.

Data structures stay simple.
The call graph maps each site to its caller, callee, argument list, and hotness.
A priority queue orders sites by score.
An incremental size model avoids remeasuring whole functions after each splice.
Attribute and alias information must survive the splice: memory effects, nonnull facts, alignment, and exception behavior travel with the copied instructions.
Debug information needs special care so stepped lines map back to the callee source even though the code now lives in the caller.
A common implementation bug is stale profiling counters after duplication; counters must be scaled or split, not copied verbatim.

Interaction with the rest of the compiler decides the final quality.
Inline before scalar replacement, loop rotation, and vectorization if those passes need the larger window.
Run dead-code elimination and constant folding immediately after each inline wave so scores for the next wave see the cleaned code.
Keep an eye on register pressure and instruction-cache footprint because an aggressive inline that helps microbenchmarks can hurt large programs.
Measure with profiles rather than assuming.
When in doubt, the maintainable default is a modest budget, eager cleanup, and specialization only for keys that simplify the clone measurably.
Good treatments of these engineering tradeoffs appear in the standard optimization texts, see~\cite{appel-1998-modern} and see~\cite{morgan-1998-optimizing}.
A further angle on preserving correctness while transforming code aggressively is given in~\cite{safonov-2010-trustworthy}.
The language rules that decide whether removing a call changes observable behavior are stated in~\cite{wg14-2011-c11}.

\section{Exercises}\label{sec:ch23-exercises}

\begin{enumerate}
\item Inline the two calls in Section~\ref{sec:ch23-worked} when the first call passes \code{s = 1} instead of \code{s = 2}. \label{ex:ch23-fold-identity} Show each renamed instruction, fold all identities, and count the surviving arithmetic operations.
\item A callee has 40 instructions and three call sites. \label{ex:ch23-budget} Two sites pass a constant flag that deletes 25 instructions after simplification, while the third passes a variable flag. Decide which sites to inline under a 50-instruction growth budget and justify the rejected site.
\item A virtual call site observes receiver types $A$, $A$, $B$, $A$ in a profile. \label{ex:ch23-spec-guard} Design a specialized clone for the hot type, draw its guard and fallback, and explain when the guard pays for itself.
\end{enumerate}

\section{Further Reading}\label{sec:ch23-reading}

Good coverage of inline costing, interprocedural context, and backend interaction appears in~\cite{morgan-1998-optimizing}.
A complementary functional-language perspective on inlining and closure handling appears in~\cite{appel-1998-modern}.
Trust and validation concerns for aggressive interprocedural rewrites are discussed in~\cite{safonov-2010-trustworthy}.
Language-level facts about function calls, argument passing, and observable behavior used to justify legality appear in~\cite{wg14-2011-c11}.
